\documentclass[12pt]{article}
\usepackage[margin=0.75in]{geometry}
\usepackage{fontspec}
\setmainfont{Latin Modern Roman}
\usepackage{microtype}
\usepackage{multicol}
\usepackage{fancyhdr}
\usepackage{titlesec}
\usepackage{enumitem}
\usepackage{xcolor}
\usepackage[hidelinks]{hyperref}
\setlength{\columnsep}{0.28in}
\setlength{\parindent}{1.1em}
\setlength{\parskip}{0.32em}
\setlength{\headheight}{15pt}
\titleformat{\section}{\large\bfseries\scshape}{}{0pt}{}
\titleformat{\subsection}{\normalsize\bfseries}{}{0pt}{}
\pagestyle{fancy}
\fancyhf{}
\fancyfoot[C]{\thepage}
\renewcommand{\headrulewidth}{0.4pt}
\newcommand{\focusbox}[1]{\par\vspace{0.4em}\noindent\begingroup\setlength{\fboxsep}{6pt}\colorbox{gray!12}{\begin{minipage}{\dimexpr\linewidth-2\fboxsep\relax}#1\end{minipage}}\endgroup\par\vspace{0.5em}}

\fancyhead[L]{Whately: Self-Deception}
\fancyhead[R]{Student Reading}
\begin{document}
\begin{center}
{\LARGE\bfseries How Characters Reason Badly About Themselves}\par\vspace{0.25em}
{\large\scshape Week 17: Literary Self-Deception}\par\vspace{0.5em}
{Applying Richard Whately's toolkit to private justification}\par\vspace{0.6em}
\rule{0.82\textwidth}{0.4pt}
\end{center}

\section*{Central Question}
How do people use reasoning to hide from themselves?

\section*{Vocabulary Preview}
\begin{multicols}{2}
\begin{itemize}[leftmargin=*, itemsep=0.18em]
\item \textbf{Self-deception}: persuading yourself of what you already want to believe.
\item \textbf{Motivated reasoning}: letting desire choose the conclusion first.
\item \textbf{Rationalization}: inventing reasons after a choice is already made.
\item \textbf{Hidden assumption}: an unstated premise the argument needs.
\item \textbf{Private argument}: the case a person makes inside a soliloquy, diary, or self-talk.
\item \textbf{Public cover}: the cleaner version offered to others.
\item \textbf{Tragic error}: a serious misjudgment that grows from defective reasoning.
\item \textbf{Tool stack}: several Whately habits used together on one case.
\item \textbf{Desire-first conclusion}: the end the speaker wants before the reasons appear.
\item \textbf{Self-deception autopsy}: map desire, claim, tools used, flaw, and repaired judgment.
\end{itemize}
\end{multicols}

\section*{Introduction}
Logic is not only for public debate. People also reason in private---to justify a choice, quiet a conscience, protect pride, or make a hard action feel necessary. Week 17 gathers the course toolkit and turns it inward: how do characters, and how do we, reason badly about ourselves?

Whately's tools remain the same: find the conclusion, test the reason, surface the hidden assumption, fix unstable terms, check form and fact, notice fallacies, weigh evidence, test sources, and separate persuasion from proof. The new pressure is motive. Desire can recruit logic as a servant.

\focusbox{\textbf{Source note:} This week applies Whately's full toolkit to self-justifying reasoning rather than introducing a new formal apparatus. Classroom and literary examples are original applications. The aim is judgment, not psychological diagnosis of students.}

\begin{multicols}{2}
\section*{Reading}

\subsection*{1. Desire Can Hire an Argument}
A person may begin with a wish: to keep a title, win an election, avoid shame, punish an enemy, or look brave. Then reasons appear. Some reasons may be true. The danger is that desire chose the destination first, and the argument is only a road built afterward.

Motivated reasoning is not always conscious lying. It often feels like honest thought. That is why the tools matter.

\subsection*{2. Private Arguments and Public Covers}
Characters often keep two versions of a case. Privately, they may admit fear, envy, or ambition. Publicly, they offer necessity, honor, safety, or the common good. A self-deception autopsy asks whether the public cover is supported independently or only restates the private wish in cleaner language.

\subsection*{3. Familiar Flaws, Private Setting}
Self-deception often reuses skills from earlier weeks:

\begin{itemize}[leftmargin=*]
\item \textbf{Hidden assumption:} ``If I want this, it must be necessary.''
\item \textbf{Equivocation:} shifting \textit{safe}, \textit{man}, \textit{honor}, or \textit{necessary}.
\item \textbf{Circular reasoning:} defining the act as justified because I am justified in doing it.
\item \textbf{Missing the point:} proving I am mistreated instead of proving I may strike.
\item \textbf{Weak analogy or thin examples:} one insult ``proves'' a whole destiny.
\item \textbf{Untrusted sources:} trusting a voice because it flatters the plan.
\item \textbf{Rhetoric over proof:} using shame or grandeur to settle a moral question.
\end{itemize}

\subsection*{4. How to Run a Self-Deception Autopsy}
Use this map on a speech, soliloquy, or everyday justification:

\begin{enumerate}[leftmargin=*,itemsep=0.12em]
\item Name the speaker's desire or feared loss.
\item State the conclusion the speaker wants to reach.
\item List the stated reasons.
\item Surface the hidden assumption.
\item Identify at least one Whately tool that exposes a flaw (term shift, circular support, weak evidence, rhetoric without proof, and so on).
\item Say what would still need discovery, evidence, or moral judgment.
\item Offer a repaired, narrower claim the evidence could support.
\end{enumerate}

\subsection*{5. Repair Without Cheap Cynicism}
The goal is not to say every self-justification is fake. Some private reasoning is honest and hard. The goal is to ask whether the argument earns its conclusion. A repaired claim is often smaller: ``I am afraid,'' ``I am ambitious,'' ``I am angry,'' or ``I see one real danger'' may be truer than ``I am forced'' or ``fate requires this.''

\subsection*{6. Why Literature Matters Here}
Tragedy is a laboratory for self-deception because characters talk themselves into irreversible acts. In \textit{Macbeth}, ambition, manhood pressure, and prophetic security language cooperate. In \textit{Julius Caesar}, Brutus dresses political fear as necessity for Rome. The plays do not replace logic class; they give dense cases where several tools must work together.

\subsection*{7. From Character to Student Habit}
The transfer is practical. Before a rushed decision, a harsh post, a portfolio excuse, or a public stand, ask: What do I want this argument to do for me? Which tool would a fair critic use first? What narrower claim can I defend without hiding?

\section*{Reading Focus}
\begin{enumerate}[leftmargin=*]
\item What is self-deception in this week's sense?
\item How does motivated reasoning differ from open lying?
\item What is the difference between a private argument and a public cover?
\item List four earlier Whately tools that often appear in self-justification.
\item Copy the seven steps of the self-deception autopsy.
\item Why can a repaired claim be narrower and still stronger?
\item Name one Shakespeare speaker who may reason badly about himself or herself.
\end{enumerate}

\section*{Handout Summary}
Write 5--7 sentences explaining how people use reasoning to hide from themselves. Your summary must include the words \textit{desire}, \textit{hidden assumption}, \textit{rationalization}, \textit{tool stack}, and \textit{repaired claim}.
\end{multicols}
\end{document}
